\chapter{プログラミングEgisonとは}

本章では，Egisonがどのような言語であるのか紹介する．
本章の内容は，いま完全に理解する必要はない．
Egisonとそれが提唱しているパターンマッチ指向プログラミングがどのような点で優れているのかその大体のイメージをつかんでもらいたい．

\section{プログラミング言語全体の中でのEgisonの位置づけ}

数多あるプログラミング言語のなかでのEgisonの著しい特徴は，世界で初めて自身に実装されたアルゴリズムの記述を直感的にするための機能を有しているところである．
アルゴリズムの記述を本質的に簡潔にするプログラミング言語のアイデアは，歴史を紐解いてみても，そう多くはない．
そういった意味でEgisonは数多にあるプログラミング言語のなかでも希少な言語であるだろう．

プログラミング言語の機能には大きく分けて2種類ある．
コンピューターを含むさまざまな機械の制御を簡潔に記述するための機能と，数学的なアルゴリズムを簡潔に記述するための機能である．
前者の機能は，たとえば，オペレーティング・システムやウェブ・ブラウザなど我々が日常的に使用しているさまざまなソフトウェアを実装する上で重要な機能である．
後者の機能は，\todo{}．
Egisonは後者の機能の拡充に注力してつくられている．

アルゴリズムを簡潔に記述することを目指しているプログラミング言語の流派の主流は，関数型プログラミング言語である．
関数型プログラミングでは，コンピューターがプログラムを実行するのかはプログラマはあまり意識せず，アルゴリズムを数学的にそのまま記述することを目指す．
関数型プログラミングでは，コンピューターが実行できる形にプログラムを書き直す役割は主にコンパイラが果たす．

関数型プログラミング言語によるプログラムの記述の大きな特徴は，関数を他の組み込みデータ型と同じようにプログラマが扱えることである．
具体的にいうと，関数を別の関数の引数として渡したり，関数を返す関数を定義することができる．
そのおかげで，関数型プログラミング言語以外の言語ではモジュール化がむずかしい処理がモジュール化できることが多々ある．

しかし，関数型プログラミング言語でも，頭のなかのアルゴリズムのイメージを，関数型プログラムとして記述できるように翻訳する必要がある場合もある．
\todo{}

\section{パターンマッチ指向プログラミングとは}

本節では，いくつかの例をみることによって，パターンマッチ指向\index{ぱたーんまっちしこう@パターンマッチ指向}プログラミングとはどのようなプログラミング・パラダイムであるのかそのイメージを紹介する．

パターンマッチ指向プログラミングによってプログラムの記述が直感的になっていることのわかりやすい例に\lstinline{intersect}関数の実装がある．
\lstinline{intersect}は2つのリストを引数にとり，それらのリストに共通する要素のリストを返す関数である．
Egisonを使って引数のリストをそれぞれ多重集合としてパターンマッチすると，2つのリストの共通要素にマッチするパターンを記述することにより，\lstinline{intersect}を記述できる．\footnote{プログラムの読み方は次章以降で紹介する．}

\begin{lstlisting}[language=egison]
intersect xs ys :=
  matchAllDFS (xs, ys) as (multiset eq, multiset eq) with
  | ($x :: _, #x :: _) -> x
\end{lstlisting}

対して，Egisonのパターンマッチを使わずに関数型プログラミングのスタイルでHaskellで記述すると，リストに対する関数を組み合わせて共通要素を抜き出すための方法を記述する必要がある．

\begin{lstlisting}[language=egison]
intersect xs ys =  filter (\x -> any (== x) ys) xs
\end{lstlisting}

このプログラムは，リスト内包表記を使って以下のように書き直すこともできる．

\begin{lstlisting}[language=egison]
intersect xs ys =  [x | x <- xs, any (== x) ys]
\end{lstlisting}

Egisonのパターンマッチを使ったプログラムは，2つのリストの共通要素にマッチするパターンを記述しているだけであるのに対し，関数型プログラミングでは，どうやってその共通要素を取り出すかその方法をプログラマが考えて記述する．
前者のように「何を計算したいか（what to do）」を記述するプログラミングスタイルは宣言型プログラミング，後者のように「どうのように計算するか（how to do）」を記述するプログラミングスタイルは手続き型プログラミングと呼ばれている．
「何を計算したいか」から「どのように計算するか」が明らかである場合は，「何を計算したいか」を記述する宣言型プログラミングのほうがプログラムが読みやすく簡潔になる．
\lstinline{intersect}の例の場合は，Egisonが多重集合のパターンマッチアルゴリズムをモジュール化できるおかげで，宣言的なプログラミングが可能になっている．

少し毛色の違う例として，\lstinline{concat}関数をパターンマッチ指向プログラミングで定義する．
リストのリストの要素にマッチするパターンを記述することにより，\lstinline{concat}を定義できる．

\begin{lstlisting}[language=egison]
concat xss :=
  matchAllDFS xss as list (list something) with
  | _ ++ (_ ++ $x :: _) :: _ -> x

concat [[1,2],[3],[4,5]]
-- [1,2,3,4,5]
\end{lstlisting}

さらにおもしろいパターンマッチ指向プログラミングの例として，\lstinline{unique}関数を定義する．
後方に自身と同じ値が含まれない要素にマッチするパターンを記述することにより\lstinline{unique}を定義できる．

\begin{lstlisting}[language=egison]
unique xs :=
  matchAllDFS xs as list eq with
  | _ ++ $x :: !(_ ++ #x :: _) -> x

unique [1,2,3,2,4]
-- [1,3,2,4]  
\end{lstlisting}

次章以降で，上記のプログラムで使われているEgisonの構文や，機能，プログラミングテクニックの解説をしていく．

\section{本書の構成}

本書は可能な限りコンパクトに，パターンマッチ指向プログラミングとそのためのEgisonの機能の解説をまとめることを目指した．
そのため，本書は前提知識として，関数型プログラミング（Lisp系言語や，OCaml，Haskellを使ったプログラミング）の知識を仮定している．
既存の関数型言語にも存在する機能については，独立した節を設けず，登場したときに軽く解説するにとどめた．

第\ref{quick}章では，パターンマッチに関係するEgisonの構文を一通り紹介する．
第\ref{pmo}章では，Egisonのパターンマッチを活かしたプログラミングスタイルであるパターンマッチ指向プログラミングとはなにか紹介する．
第\ref{pmo-challange}章では，より実践的なパターンマッチ指向プログラミングの例を紹介する．
第\ref{matcher}章では，新しいマッチャーの定義の方法を説明する．
第\ref{implementation}章では，Egisonのパターンマッチの実装について紹介する．
